% !TeX encoding = UTF-8
% Español de España. Cargar un único arquetipo desde el archivo principal.
\newcommand{\cvsummary}{Ingeniero mecánico y mecatrónico por NYU Abu Dhabi, con experiencia en sistemas mecánicos en los Emiratos Árabes Unidos y Europa. Experiencia práctica en mejora de la eficiencia energética industrial, desarrollo de prototipos de energía eólica y apoyo a proyectos de infraestructuras hidráulicas.}


% --- COMPETENCIAS ---
\newcommand{\skills}{
    \section*{COMPETENCIAS TÉCNICAS}
    {\small
    \textbf{Energía y sostenibilidad:} Eólica marina flotante (estabilidad hidrostática de plataformas semisumergibles y validación en canal de oleaje), aerodinámica de aerogeneradores de eje vertical (VAWT; patente en trámite), sistemas térmicos (prototipos TEC y de climatización), corrección del factor de potencia y análisis de subestaciones de alta tensión, análisis del mercado de captura de carbono (CCUS), ecodiseño y metrología de medición de carbono.

    \vspace{1pt}
    \textbf{Simulación y análisis:} ANSYS (Mechanical, Fluent, AQWA), COMSOL Multiphysics, análisis por elementos finitos (FEA), dinámica de fluidos computacional (CFD), análisis térmico y multifísico, modelado dinámico con MATLAB/Simulink.

    \vspace{1pt}
    \textbf{CAD y diseño:} SolidWorks (CSWA), CATIA V5, Siemens NX, CREO, MoldFlow, Fusion 360, AutoCAD Electrical, diseño paramétrico 3D, análisis estructural y cinemático, tolerancias geométricas y dimensionales (GD\&T).

    \vspace{1pt}
    \textbf{Datos y SIG:} PostgreSQL/PostGIS, Python, MATLAB, procesos de tratamiento de datos en SQL, integración de datos geoespaciales (SIG de infraestructuras hidráulicas) y visualización de datos.

    \vspace{1pt}
    \textbf{Control e instrumentación:} Control PID, fusión sensorial con IMU, Arduino, PLC, sensores e instrumentación, integración de hardware y software.

    \vspace{1pt}
    \textbf{Gestión de proyectos:} MS Project, Jira, Notion, coordinación multidisciplinar, informes técnicos para clientes, Lean Six Sigma White Belt.

    \vspace{1pt}
    \textbf{Idiomas:} Español nativo, inglés C2, francés C2, italiano A1 y alemán A1.
    }
}



% --- PROYECTOS ---

\newcommand{\missiletitle}{Simulador de dinámica de interceptación de misiles}
\newcommand{\missilelocation}{Besançon, Francia}
\newcommand{\missiledescOne}{Desarrollé en MATLAB una simulación 3D de interceptación de misiles mediante guiado por navegación proporcional, integración numérica y visualización de trayectorias en tiempo real con detección de impacto.}

\newcommand{\vawttitle}{VAWT -- Aprovechamiento de energía eólica urbana}
\newcommand{\vawtdates}{Otoño de 2024 - Primavera de 2025}
\newcommand{\vawtlocation}{Abu Dabi, EAU}
\newcommand{\vawtdescOne}{Desarrollé un sistema de aprovechamiento de energía eólica urbana con solicitud de patente presentada para captar el viento generado por el tráfico, mediante el diseño de un aerogenerador de eje vertical dimensionado para las ráfagas junto a la calzada.}
\newcommand{\vawtdescTwo}{Validé el rendimiento antes de la fabricación mediante simulaciones CFD en ANSYS Fluent contrastadas con datos de viento medidos, analizando el comportamiento del rotor en regímenes realistas de viento generado por el tráfico.}

\newcommand{\fsaetitle}{Formula SAE -- Suspensión y chasis de fibra de carbono}
\newcommand{\fsaedates}{Otoño de 2024 -- Actualidad}
\newcommand{\fsaelocation}{Abu Dabi, EAU}
\newcommand{\fsaedescOne}{Desarrollé un conjunto de suspensión y chasis de fibra de carbono, con prestaciones estructurales verificadas mediante análisis por elementos finitos, diseñando y simulando el ensamblaje en SolidWorks y ANSYS.}
\newcommand{\fsaedescTwo}{Ajusté la dinámica del vehículo mediante estudios cinemáticos en SolidWorks, evaluando las curvas de caída, balanceo y relación de movimiento a lo largo del recorrido de la suspensión.}

\newcommand{\tactitle}{Mininevera -- Sistema termoeléctrico TEC de dos etapas}
\newcommand{\tacdates}{Primavera de 2025}
\newcommand{\taclocation}{Abu Dabi, EAU}
\newcommand{\tacdescOne}{Eliminé por completo la carga de refrigerante mediante un diseño de refrigeración Peltier de dos etapas en estado sólido, sin fluido de trabajo, basado en módulos TEC.}
\newcommand{\tacdescTwo}{Gestioné el calor residual mediante el diseño y la impresión 3D de los disipadores y de toda la estructura de soporte, supervisando su rendimiento térmico.}

\newcommand{\orbitaltitle}{OrbitalGrid -- Sistema de información geográfica espacial}


\newcommand{\orbitaldates}{06/2024 - Actualidad}
\newcommand{\orbitallocation}{Abu Dabi, EAU}
\newcommand{\orbitaldescOne}{Desarrollé una base de datos espacial para cartografía aeroespacial, con capas de escaneado 3D y datos satelitales consultables en PostgreSQL, diseñando el esquema y los procesos de carga de datos.}
\newcommand{\orbitaldescTwo}{Habilité el análisis entre distintas fuentes mediante cruces entre geometría escaneada e imágenes satelitales, normalizando los sistemas de coordenadas y la indexación espacial.}

\newcommand{\softrobtitle}{Exoesqueleto robótico para rehabilitación tras un ictus}
\newcommand{\softrobdates}{04/2024 - 05/2024}
\newcommand{\softroblocation}{Abu Dabi, EAU}
\newcommand{\softrobdescOne}{Restablecí la flexión de los dedos mediante actuadores para pacientes en rehabilitación, evaluando el rango de movimiento articular de un exoesqueleto inspirado en el origami y accionado por hilos de aleación con memoria de forma (SMA).}
\newcommand{\softrobdescTwo}{Reproduje la flexibilidad del tejido humano mediante la combinación de hilos SMA y polímero Dragon Skin en lugar de articulaciones rígidas, evaluando la transmisión de fuerzas a través de una interfaz blanda.}

\newcommand{\exoplanettitle}{Simulador de tránsitos de exoplanetas}
\newcommand{\exoplanetdates}{Primavera de 2025}
\newcommand{\exoplanetlocation}{Abu Dabi, EAU}
\newcommand{\exoplanetdescOne}{Simulé tránsitos de exoplanetas mediante la implementación de mecánica orbital y física estelar en Python, obteniendo curvas de luz con la profundidad y la duración previstas.}
\newcommand{\exoplanetdescTwo}{Generalicé el modelo para distintas configuraciones mediante su parametrización, en lugar de programar casos fijos, y ejecuté simulaciones con diferentes parámetros de estrella y planeta.}

\newcommand{\windTurbinetitle}{Aerogenerador marino flotante -- Estabilización activa}
\newcommand{\windTurbinedates}{05/2026 - 06/2026}
\newcommand{\windTurbinelocation}{Besançon, Francia}
\newcommand{\windTurbinedescOne}{Garanticé la estabilidad pasiva en cabeceo de una plataforma semisumergible a escala 1:1000, con altura metacéntrica GM = 5,38 cm y rigidez de restitución K = 0,0607 N$\cdot$m/rad, mediante el cálculo hidrostático de una geometría de tres columnas tipo WindFloat en SolidWorks.}
\newcommand{\windTurbinedescTwo}{Estabilicé el movimiento de una plataforma semisumergible en un canal de oleaje mediante control en lazo cerrado del cabeceo y el balanceo de un prototipo instrumentado, integrando control de dos ejes con anti-windup y gestión de zona muerta y saturación.}

\newcommand{\carbonmetertitle}{CARBONMETRE -- Ecodiseño y medición de carbono}
\newcommand{\carbonmeterdates}{09/2025 - 12/2025}
\newcommand{\carbonmeterlocation}{Besançon, Francia}
\newcommand{\carbonmeterdescOne}{Desarrollé un método de medición de carbono para una pieza fabricada, con un procedimiento trazable desde las etapas de fabricación hasta el valor comunicado, combinando datos de fabricación industrial y prácticas metrológicas.}
\newcommand{\carbonmeterdescTwo}{Relacioné las decisiones de fabricación con su impacto medido mediante la comparación del carbono incorporado de distintas rutas de fabricación, estructurando los datos de materiales y procesos para facilitar una comparación directa.}

\newcommand{\hvactitle}{Acondicionamiento térmico de una habitación de residencia -- Refrigeración pasiva y evaporativa}
\newcommand{\hvacdates}{06/2026 - 06/2026}
\newcommand{\hvaclocation}{Besançon, Francia}
\newcommand{\hvacdescOne}{Reduje el sobrecalentamiento de una habitación de residencia en horas punta, con una temperatura 3 \textdegree{}C inferior a la exterior entre las 14:00 y las 21:00, mediante una pantalla exterior de lámina reflectante, EPS de 1 cm y cartón ondulado.}
\newcommand{\hvacdescTwo}{Bloqueé la radiación solar directa sobre una ventana de vidrio simple de 9 m$^2$ mediante capas exteriores reflectantes y aislantes, verificando la reducción de la temperatura superficial en las horas de máxima insolación.}

\newcommand{\controlBoxtitle}{Instalación eléctrica integral -- Sistemas para eventos de mediana escala}
\newcommand{\controlBoxdates}{06/2023 - 08/2023}
\newcommand{\controlBoxlocation}{Ciudad de México, México}
\newcommand{\controlBoxdescOne}{Reduje un 30\% el tiempo de puesta en marcha eléctrica de eventos, medido en horas de montaje antes de la apertura al público, mediante el diseño de un cuadro de control con PLC de 12 entradas.}
\newcommand{\controlBoxdescTwo}{Desarrollé una instalación completa de mediana escala, con un diseño de distribución listo para su puesta en servicio, dimensionando circuitos para cargas audiovisuales y alimentación del escenario.}

\newcommand{\hyperlooptitle}{NYC Hyperloop -- Vehículo a escala reducida}
\newcommand{\hyperloopdates}{01/2024 - 05/2024}
\newcommand{\hyperlooplocation}{Nueva York, EE. UU.}
\newcommand{\hyperloopdescOne}{Desarrollé un vehículo Hyperloop a escala reducida para un concepto de trayecto Nueva York--Long Island, liderando el subequipo mecánico hasta la construcción e integración del prototipo.}
\newcommand{\hyperloopdescTwo}{Verifiqué el comportamiento de levitación y aerodinámico mediante estudios acoplados electromagnéticos y de flujo en COMSOL y ANSYS, cuyos resultados orientaron las modificaciones de diseño.}

\newcommand{\dronetitle}{Cuadricóptero funcional -- Estructura ligera}
\newcommand{\dronedates}{03/2024 - 05/2024}
\newcommand{\dronelocation}{Nueva York, EE. UU.}
\newcommand{\dronedescOne}{Reduje un 40\% la masa de la estructura respecto al diseño de referencia, comparando el peso de las piezas impresas, mediante fabricación aditiva en PLA y TPU.}
\newcommand{\dronedescTwo}{Cumplí los criterios de vibración y caída mediante iteraciones geométricas en SolidWorks y ANSYS, contrastando los resultados de ensayo con los umbrales de rendimiento establecidos.}

\newcommand{\trentxwbtitle}{Rolls-Royce Trent XWB -- Modelo 3D funcional}
\newcommand{\trentxwbdates}{01/2024 - 05/2024}
\newcommand{\trentxwblocation}{Nueva York, EE. UU.}
\newcommand{\trentxwbdescOne}{Reproduje el Trent XWB como un conjunto funcional completo en SolidWorks y Fusion 360, cubriendo todos sus módulos mediante ingeniería inversa a partir de planos públicos.}
\newcommand{\trentxwbdescTwo}{Mejoré la eficiencia del diseño mediante estudios de flujo y estructurales en ANSYS sobre secciones críticas, evaluando los resultados aerodinámicos y de deformación simulados.}

\newcommand{\liostitle}{LIOS -- Asistente de voz con IA}
\newcommand{\liosdates}{12/2024 - Actualidad}
\newcommand{\lioslocation}{Abu Dabi, EAU}
\newcommand{\liosdescOne}{Desarrollé un asistente de voz funcional en Python con procesamiento del lenguaje natural e integración de API, capaz de completar tareas de uso diario sin intervención manual.}
\newcommand{\liosdescTwo}{Mantuve el contexto en conversaciones de varios turnos mediante memoria contextual, conservando la información de la sesión entre consultas.}

\newcommand{\fitpaltitle}{FitPal -- Dispositivo inteligente para entrenamiento físico}
\newcommand{\fitpaldates}{09/2021 - 12/2021}
\newcommand{\fitpallocation}{Abu Dabi, EAU}
\newcommand{\fitpaldescOne}{Codesarrollé un dispositivo inteligente de entrenamiento físico basado en M5Stack y su firmware, con información sobre la ejecución de ejercicios en tiempo real.}
\newcommand{\fitpaldescTwo}{Transmití datos de sensores a una aplicación móvil en tiempo real mediante Arduino IoT Cloud, integrando los puntos de conexión IoT.}

\newcommand{\covidtitle}{Cabina de desinfección automatizada}
\newcommand{\coviddates}{09/2020 - 02/2021}
\newcommand{\covidlocation}{Ciudad de México, México}
\newcommand{\coviddescOne}{Desarrollé un prototipo funcional de cabina de desinfección automatizada controlada por Arduino, integrando sensores, nebulización y control inalámbrico.}
\newcommand{\coviddescTwo}{Reduje la exposición del operario mediante la activación inalámbrica y sin contacto del ciclo de desinfección.}

\newcommand{\rovertitle}{Rover teleoperado -- Escaneado espacial 3D}
\newcommand{\roverdates}{09/2026 - Actualidad}
\newcommand{\roverlocation}{Gijón, España}
\newcommand{\roverdescOne}{Implementé la captura 3D completa desde una plataforma controlada a distancia, adquiriendo nubes de puntos sin operarios en la zona de escaneado mediante la integración del control embebido. Diseñé todo el chasis y el sistema mecánico del rover en SolidWorks y realicé análisis estructurales en ANSYS.}
\newcommand{\roverdescTwo}{Mantuve al operario fuera del entorno de medición mediante el control remoto completo del movimiento y la adquisición de datos, desarrollando el enlace de teleoperación y la interfaz de mando.}

\newcommand{\microsensortitle}{Microsensor piezorresistivo de presión de detonación}
\newcommand{\microsensordates}{05/2026 - 06/2026}
\newcommand{\microsensorlocation}{Besançon, Francia}
\newcommand{\microsensordescOne}{Desarrollé un diseño de sensor de presión para entornos de detonación de 10 GPa, cumpliendo simultáneamente los requisitos de supervivencia y de interfaz de 50 $\Omega$, mediante la selección de todas las capas de materiales bajo restricciones eléctricas, mecánicas y económicas acopladas.}
\newcommand{\microsensordescTwo}{Diseñé un microsensor piezorresistivo y automaticé las iteraciones de geometría y máscaras mediante código que impone restricciones de K, $\Delta R/R_0$ y resistividad; trasladé las especificaciones a una máscara litográfica y una secuencia de proceso a escala de oblea. Calificación de 20/20 en microfabricación.}

% --- EXPERIENCIA ---
\newcommand{\carbonorg}{Energy Innovation Capital}
\newcommand{\carbontitle}{Analista de mercado de captura de carbono}
\newcommand{\carbondates}{10/2023 - 01/2024}
\newcommand{\carbonlocation}{Abu Dabi, EAU}
\newcommand{\carbonbullets}{%
  \item Elaboré recomendaciones estratégicas de inversión en empresas emergentes de captura de carbono mediante una comparación homogénea de las vías CCUS a lo largo de la cadena de valor, presentando una selección de empresas a analistas sénior.
  \item Cartografié el mercado CCUS con una visión estructurada de los actores de captura, transporte y almacenamiento, contrastando bibliografía científica y anuncios de proyectos comerciales.
  \item Distinguí los avances técnicos reales de las afirmaciones comerciales mediante una comparación de la madurez tecnológica de las vías de captura, revisando los datos publicados de rendimiento y costes.
  \item Transformé resultados de ingeniería en informes para la toma de decisiones, traduciendo los hallazgos técnicos al lenguaje de mercado y resumiendo los riesgos de activos de tecnología climática para analistas.
}

\newcommand{\spacelaborg}{NYU Abu Dhabi, Centro de Ciencias Espaciales}
\newcommand{\spacelabtitle}{Ingeniero e investigador de I+D -- Plataforma de ciencias planetarias}
\newcommand{\spacelabdates}{05/2025 - 09/2025}
\newcommand{\spacelablocation}{Abu Dabi, EAU}
\newcommand{\spacelabbullets}{%
  \item Reduje un 40\% el tiempo del proceso de tratamiento de datos, evaluado de principio a fin sobre conjuntos de imágenes a escala planetaria, mediante algoritmos de compresión y reproyección asistidos por aprendizaje automático.
  \item Desarrollé y desplegué un atlas de Marte con imágenes multiespectrales reales de la misión Emirates Mars Mission HOPE, desde el diseño hasta una \href{https://romvio.github.io/mars-viewer/}{aplicación web y móvil} operativa y una versión iOS para la App Store de Apple en EAU.
  \item Mejoré la fidelidad de los mapas entre pasadas orbitales solapadas, con mayor uniformidad de brillo y continuidad de las uniones, mediante calibración de solapamientos y reproyección cilíndrica a Robinson.
  \item Fundamenté el atlas en datos físicos mediante la recogida de muestras análogas del desierto de EAU y la caracterización de la estructura y morfología de sus granos para emular el regolito marciano.
}

\newcommand{\geoorg}{Global Education, NYU Abu Dhabi}
\newcommand{\geotitle}{Coordinador de Global Education -- Estudiantes visitantes}
\newcommand{\geodates}{12/2025 - 05/2025}
\newcommand{\geolocation}{Abu Dabi, EAU}
\newcommand{\geobullets}{%
  \item Organicé programas académicos y culturales de inmersión en EAU para estudiantes internacionales visitantes, con salidas realizadas según el calendario del semestre, coordinando transporte, recintos y permisos de principio a fin.
  \item Reduje los riesgos de los desplazamientos en grupo mediante protocolos de seguridad ejecutados sin incidentes, liderando el equipo de coordinadores y asignando responsables por emplazamiento.
  \item Actué como interlocutor único para los estudiantes visitantes, resolviendo incidencias en el primer contacto mediante coordinación directa con los servicios del campus.
  \item Mantuve la coordinación entre las unidades anfitrionas y los traspasos entre departamentos, centralizando la comunicación en un canal por grupo.
}

\newcommand{\geoOneorg}{Global Education, NYU Abu Dhabi}
\newcommand{\geoOnetitle}{Coordinador de Global Education -- Estudiantes de primer curso}
\newcommand{\geoOnedates}{09/2026 - 12/2024}
\newcommand{\geoOnelocation}{Abu Dabi, EAU}
\newcommand{\geoOnebullets}{%
  \item Facilité la transición de estudiantes de primer curso mediante sesiones estructuradas de orientación al inicio del semestre, en colaboración con el equipo de Global Education.
  \item Organicé y dirigí personalmente visitas y sesiones de asesoramiento semanales para familiarizar a los estudiantes con el campus.
  \item Promoví la integración en la comunidad mediante la planificación de actividades académicas y culturales y eventos de inclusión para distintos grupos.
  \item Ofrecí apoyo continuado de orientación entre estudiantes durante el semestre mediante horarios periódicos de atención sin cita previa.
}

\newcommand{\hcorg}{Centro de Salud, NYU Abu Dhabi}
\newcommand{\hctitle}{Empleo de estudiante -- Responsable de redes sociales}
\newcommand{\hcdates}{09/2024 - 05/2025}
\newcommand{\hclocation}{Abu Dabi, EAU}
\newcommand{\hcbullets}{%
  \item Aumenté la participación y el alcance de las campañas de salud mediante contenidos de bienestar adaptados a distintos segmentos del alumnado.
  \item Mantuve una frecuencia regular de publicación conforme al calendario mensual de contenidos, elaborando internamente los textos y materiales gráficos.
  \item Reduje los tiempos de respuesta a consultas en el buzón compartido, actuando como contacto principal y derivando los casos correspondientes.
  \item Mantuve un tono uniforme en todos los canales mediante la aplicación de las directrices de comunicación clínica del centro de salud.
}

\newcommand{\alpesorg}{Alpes Technologies México}
\newcommand{\alpestitle}{Ingeniero de proyectos júnior -- SGA y sistemas eléctricos}
\newcommand{\alpesdates}{06/2024 - 09/2024}
\newcommand{\alpeslocation}{Ciudad de México, México}
\newcommand{\alpesbullets}{%
  \item Elevé el factor de potencia de los clientes al 99\%, con una reducción del 15\% en costes energéticos y del 25\% en paradas eléctricas, mediante la especificación e instalación de baterías de condensadores francesas tras analizar armónicos y picos de tensión en subestaciones de alta tensión.
  \item Sustituí un proceso de inventario completamente manual por un sistema de gestión de almacenes (SGA) propio en HTML/JS y Google Sheets, aumentando un 60\% la eficiencia operativa y de fabricación.
  \item Implementé trazabilidad por lotes para el control de calidad, cubriendo en tiempo real todos los movimientos de existencias con códigos QR, actualizaciones bidireccionales y acceso por roles en una sola herramienta.
  \item Restablecí la instalación eléctrica de un complejo hotelero afectado por un incendio mediante mediciones de campo durante la emergencia, elaborando esquemas unifilares y recomendaciones correctivas aceptados por el consejo de administración del cliente.
}

\newcommand{\epigsaorg}{EPIGSA}
\newcommand{\epigsatitle}{Asistente de ingeniería de datos -- SIG e hidráulica}
\newcommand{\epigsadates}{05/2024 - 07/2024}
\newcommand{\epigsalocation}{Ciudad de México, México}
\newcommand{\epigsabullets}{%
  \item Consolidé datos de infraestructuras municipales de tratamiento de aguas mediante la integración y el contraste de fuentes de levantamiento, generando capas SIG validadas y cargadas en PostgreSQL.
  \item Reduje un 25\% los errores de datos, medidos tras la validación automatizada, mediante rutinas de aseguramiento y control de calidad en el proceso de ingesta.
  \item Establecí un flujo de trabajo repetible y documentado para el equipo, normalizando el esquema de datos y las convenciones de nomenclatura.
  \item Proporcioné una base geométrica verificada para los proyectos mediante la conciliación de registros de fuentes contradictorias, resolviendo duplicados y discrepancias de proyección en origen.
}

\newcommand{\solarshiporg}{Solar Ship}
\newcommand{\solarshiptitle}{Ingeniero en prácticas de estrategia y diseño de movilidad aérea}
\newcommand{\solarshipdates}{05/2023 - 08/2023}
\newcommand{\solarshiplocation}{Toronto, Canadá}
\newcommand{\solarshipbullets}{%
  \item Diseñé un concepto de logística aérea totalmente eléctrica para comunidades sin conexión a la red, con rutas y cargas útiles en zonas remotas de Zambia y Tanzania, dimensionando sistemas de aeronaves eléctricas para pistas sin servicio.
  \item Integré la sostenibilidad en el concepto de operaciones mediante criterios de impacto social incluidos en la evaluación de alternativas de diseño, tratando el acceso de las comunidades como una restricción de diseño y no como un resultado.
  \item Contribuí a un programa multinacional de ingeniería mediante entregables de diseño para el conjunto del equipo, armonizando las hipótesis de los subsistemas entre disciplinas.
  \item Justifiqué la infraestructura necesaria para conexiones aéreas remotas mediante una comparación documentada con el transporte terrestre, evaluando tiempos de acceso y costes por ruta.
}

\newcommand{\crousorg}{CROUS Besançon Bourgogne Franche-Comté}
\newcommand{\croustitle}{Empleo de estudiante -- Restauración universitaria}
\newcommand{\crousdates}{10/2024 - Actualidad}
\newcommand{\crouslocation}{Besançon, Francia}
\newcommand{\crousbullets}{%
  \item Mantuve el ritmo de trabajo de cocina en horas punta, garantizando la continuidad del servicio de comidas mediante apoyo en los puestos de lavado y preparación.
  \item Cumplí los requisitos de higiene y superé los controles de seguridad alimentaria mediante la aplicación de procedimientos de limpieza y desinfección.
  \item Reduje las interrupciones del servicio mediante la reposición de suministros durante la actividad y el apoyo directo a la atención al público.
  \item Mantuve una asistencia constante durante el curso académico, adaptándome a la rotación de puestos y turnos.
}

\newcommand{\rallyorg}{Campeonato Mundial de Rally -- Guanajuato}
\newcommand{\rallytitle}{Comisario técnico oficial de la FIA}
\newcommand{\rallydates}{11/2019 - 04/2020}
\newcommand{\rallylocation}{Guanajuato, México}
\newcommand{\rallybullets}{%
  \item Garanticé el cumplimiento técnico, sin vehículos no conformes en la salida, mediante el marcaje oficial y la inspección de neumáticos conforme a procedimientos ISO 9001.
  \item Completé las inspecciones de los vehículos antes de cada sesión dentro de los plazos del fin de semana de competición, coordinándome con los comisarios responsables.
  \item Comuniqué con claridad las decisiones de inspección, aceptadas por pilotos y equipos, indicando la normativa incumplida y la corrección necesaria.
  \item Mantuve registros de marcaje trazables y auditables por participante, anotando los datos en el momento de la inspección.
}

\newcommand{\papaloteorg}{Papalote Museo del Niño}
\newcommand{\papalotetitle}{Auxiliar de atención al público}
\newcommand{\papalotedates}{05/2019 - 07/2019}
\newcommand{\papalotelocation}{Ciudad de México, México}
\newcommand{\papalotebullets}{%
  \item Expliqué exposiciones de ciencia y tecnología a niños mediante demostraciones prácticas que mantuvieron su participación durante las visitas.
  \item Mejoré la experiencia de los visitantes mediante propuestas y pruebas con el equipo de operaciones, con cambios en sala implantados tras la revisión del personal.
  \item Atendí a un público diverso, adaptando las explicaciones a distintas edades y necesidades.
  \item Mantuve un funcionamiento seguro de las exposiciones mediante supervisión sin incidentes y aplicación de las normas de acceso a los módulos interactivos.
}
